\begin{comment}
=== AI-INSTRUCTIONS (NEVER DELETE OR ALTER THIS BLOCK) ===
1. VERBATIM RETENTION: Always include this exact comment block at the top of your output.
2. LHS STYLE: Explanations on the Left-Hand Side must be colloquial, punchy, and use "footnotesize".
3. DRY/KISS: Always tie the code back to SE principles (like DRY, KISS, YAGNI, SRP) where applicable.
4. SPATIAL LIMITS: Keep the minted code to ~40 lines maximum.
5. REVIEW NOTES:
   - [Tim note2]: in SE notes, point how how mentioned TLAs are used in other, much larger systems.
   - [Tim's Note]: Make sure to emphasize why we aren't using Pandas (YAGNI).
   - [Tim's Note]: Keep the magic characters (+, -, !) explanation prominent.
   - [T note]: early on say ezr is a data-lite, model-lite, code-lite approach to AI.
              data-lite: not assuming terabyres of data. many situation si se when we have just few dozen reliable labels, at most. 
              see [^aa], below.
              model-lite: instead of using (say) 100s of decision trees in an ensemble, like SMAC3 [^smacs3], 
               or trillions of models, like LLMs
              code-lite: it tiny data and models, should be able to write tiny scripts to ahndle it all, especially since we've had
               decades of exeorence in learning hw to cod these things better and what 
               shortcuts to take. connect to KISS. YAGNI

   - [T note]: got to say early on that this doe not handle all ai.
               no, this does not handle all ai tasks but it does ande a lot  (including software configuration, performance tuning, product line engineering,
               project health forecasting, defect prediction, software testing, software process and cost estimation, and
               cross-domain generalization datasets).
               so someties we need things as complex as possible but some often we tend to complexitt
               via cogitive bias rather than after comparing simpler to complex

   - [T note]: use of clone.
   - [T note]: "Data" extrardoanirly useful . storing data from disk. clustering: one new Data be clsuter.
               decision tree: one new Data
               per tree node.
               Bayesclassifier: store distributions in differernt classes. Simualted annearling: store distributions of backgrund data
               so we now what ia sensible size of mutation.
   - [T note}: need more on the iter idioms, next. iter.
               (generators)
   - [T note]:  there are many constructions of first code sample that would confuse/mislead newbies. e.g.
               comprehensions. comprejensions with 
               conditionals.
               type
               all such idoms need to be teased out.
               maybe takes more than 1 page

   - [T note]:  centroids, cached. reused.
               only recacled if Data contents changed.
   - [T note]:  mention this is active learninger (best results when learners pick their own data. 
                by focusing on most informative examples, we can dodge noise and irrelvancies and superflous signnals and learn from
                vastly less data.
     

[^aa]: find 𝑋 Find associated 𝑌 values It can be very quick to ... It is
a much slower task to...  
Mine GitHub to find all the distributions
of code size, num- ber of dependencies per function, etc.  Discover
(a)how much that software could be sold on the market or (b) what
is the time required to build this kind of software Count the number
of classes in a system. Negotiate with an organization permission
to find how much human effort was required to build and maintain
that code.  List design options; e.g. 20 binary choices is 220 >
1, 000, 000 options.  Check all those options with a group of human
stakehold- ers.  List the configuration parameters for some piece
of soft- ware.  Generate 
a separate executable for each one of those
pa- rameter settings, then run those executable through some test
suite.
List the controls of a data miners used in software analytics
(e.g. how many neighbors to use in a k-th nearest neighbor classifier).
Run a grid search looking for the best settings for some local data.
Generate test case inputs (e.g.) using some grammar-based fuzzing.
Run all all those tests.
It can be even slower for a human to check
through all those results looking for anomalous behavior.
[^smac3]: @article{JMLR:v23:21-0888,
  author  = {Marius Lindauer and Katharina Eggensperger and Matthias
  Feurer and André Biedenkapp and Difan Deng and Carolin Benjamins
  and Tim Ruhkopf and René Sass and Frank Hutter},
    title   = {SMAC3: A Versatile Bayesian Optimization Package for
    Hyperparameter Optimization},
      journal = {Journal of Machine Learning Research},
    year    = {2022},
      volume  = {23},
        number  = {54},
          pages   = {1--9},
        url     = {http://jmlr.org/papers/v23/21-0888.html}
        }

==========================================================
\end{comment}

% ==========================================
% SLIDE 1: Intent & Motivation
% ==========================================
\begin{frame}{1. Intent \& Motivation: The "Lite" Approach}
  \vspace{0.3cm}
  \textbf{The Goal:} \texttt{ezr} is a \textbf{data-lite, model-lite, code-lite} approach to AI.
  
  \vspace{0.2cm}
  \begin{itemize}
      \item \textbf{Data-lite:} We don't assume terabytes of data. In SE, we often only have a few dozen reliable labels.
      \item \textbf{Model-lite:} We aren't building ensembles of 100s of decision trees (like SMAC3), nor are we building trillion-parameter LLMs.
      \item \textbf{Code-lite:} Tiny data + tiny models = tiny scripts. Heavy libraries like Pandas violate our core tenets (\textbf{YAGNI}).
  \end{itemize}

  \vspace{0.4cm}
  \textbf{Does this handle all AI?} \\
  No. But it handles a \textit{massive} slice of SE tasks: configuration, tuning, defect prediction, project forecasting, and test generation. We often default to massive complexity via cognitive bias, rather than after genuinely comparing simple vs. complex!
\end{frame}

% ==========================================
% SLIDE 2: Core Architecture (Code)
% ==========================================
\begin{frame}[fragile]{2. Core Architecture: The "Data" Workhorse}
\begin{columns}[T]
  \begin{column}{0.4\textwidth}
    \footnotesize 
    \textbf{Lines 145-163:}\\
    The \texttt{Data} structure is extraordinarily useful. It's the core atom of our system.

    \vspace{0.2cm}
    \rc{1} \textbf{Boss Class:} \texttt{Data} grabs headers and streams the rest.
    
    \vspace{0.15cm}
    \rc{2} \textbf{The Clone Wars:} \texttt{clone} copies the structure. Every cluster, tree node, and Bayes class is just another \texttt{Data} object!

    \vspace{0.15cm}
    \rc{3} \textbf{Magic Sorting:} Headers with \texttt{!} (maximize) or \texttt{-} (minimize) become dependent (\texttt{ys}); others independent (\texttt{xs}).
  \end{column}
  
  \begin{column}{0.6\textwidth}
\begin{minted}{python}
# ---- 2. Data (Tables) ----
class Data: |\rc{1}|
  def __init__(i, src: Iterable = None):
    src = iter(src or {})
    i.rows, i.cols, i._centroid = [], Cols(next(src)), None
    adds(src, i)

def clone(d: Data, rs: list = None) -> Data: |\rc{2}|
  return adds(rs or [], Data([d.cols.names]))

def mids(d: Data) -> list[Atom]:
  d._centroid = d._centroid or [mid(c) for c in d.cols.all]
  return d._centroid

class Cols:
  def __init__(i, names: list[str]):
    i.names, i.all = names, [Col(txt, j) for j, txt in enumerate(names)]
    i.klass = next((c for c in i.all if c.txt[-1] == "!"), None)
    
    # Separation by header notation
    i.xs = [c for c in i.all if c.txt[-1] not in "+-!X"] |\rc{3}|
    i.ys = [c for c in i.all if c.txt[-1] in "+-!"]
\end{minted}
  \end{column}
\end{columns}
\end{frame}

% ==========================================
% SLIDE 3: Scripting Notes A
% ==========================================
\begin{frame}[fragile]{3. Scripting Notes: Peeling the Iterator}
  \vspace{0.2cm}
  Our first code block uses constructions that can mislead newbies. Let's tease them out.

  \vspace{0.3cm}
  \textbf{Idiom 1: Safe Iterators \& The \texttt{next()} peel}
  \begin{minted}[fontsize=\footnotesize]{python}
  src = iter(src or {})
  i.rows, i.cols = [], Cols(next(src))
  \end{minted}
  We don't load the whole file. We wrap it in a generator (\texttt{iter}). Calling \texttt{next()} cleanly "peels" off the first row (the headers) to build our columns. The rest remains lazily waiting to be streamed (\textbf{KISS} memory management).
  
  \vspace{0.3cm}
  \textbf{Idiom 2: Lazy Generator Extraction}
  \begin{minted}[fontsize=\footnotesize]{python}
  i.klass = next((c for c in i.all if c.txt[-1] == "!"), None)
  \end{minted}
  Instead of a bulky \texttt{for}-loop, this builds a lazy generator expression \textit{inside} \texttt{next(...)}. It stops evaluating at the very first match, or safely defaults to \texttt{None}.
\end{frame}

% ==========================================
% SLIDE 4: Scripting Notes B
% ==========================================
\begin{frame}[fragile]{4. Scripting Notes: Comprehensions \& Caches}
  
  \vspace{0.2cm}
  \textbf{Idiom 3: Conditional List Comprehensions}
  \begin{minted}[fontsize=\footnotesize]{python}
  i.xs = [c for c in i.all if c.txt[-1] not in "+-!X"]
  \end{minted}
  List comprehensions with inline conditionals filter data at C-speed. It checks the last character (\texttt{txt[-1]}) to instantly route variables. Pure \textbf{DRY} logic.

  \vspace{0.4cm}
  \textbf{Idiom 4: Lazy Caching}
  \begin{minted}[fontsize=\footnotesize]{python}
  def mids(d: Data) -> list[Atom]:
    d._centroid = d._centroid or [mid(c) for c in d.cols.all]
  \end{minted}
  We use \texttt{or} for short-circuit evaluation. Centroids are cached and reused! They are \textit{only} recalculated if \texttt{d.\_centroid} is cleared (which happens exactly when \texttt{Data} contents change).
\end{frame}

% ==========================================
% SLIDE 5: SE Notes
% ==========================================
\begin{frame}{5. SE Notes: Software Engineering Principles}
  \vspace{0.2cm}
  This codebase thrives by strictly adhering to core SE principles (TLAs) that also govern massive systems:

  \vspace{0.3cm}
  \begin{block}{\textbf{DRY} (Don't Repeat Yourself)}
  The \texttt{Data} object is reused relentlessly via \texttt{clone} for trees, clusters, and Bayes. \textit{In massive systems like Linux or Kubernetes, adhering strictly to DRY is exactly what prevents cascading bug failures across million-line codebases.}
  \end{block}

  \begin{block}{\textbf{KISS} \& \textbf{YAGNI} (Keep It Simple / You Ain't Gonna Need It)}
  We didn't import Pandas. We just needed an iterable to split \texttt{Xs} from \texttt{Ys}. \textit{Even tech giants use YAGNI to prevent feature bloat. Our code-lite approach proves that SE experience allows smart shortcuts over defaulting to bloated enterprise frameworks.}
  \end{block}

  \begin{block}{\textbf{SRP} (Single Responsibility Principle)}
  \texttt{Data} strictly stores rows. \texttt{Cols} strictly parses schemas. They do not overlap. \textit{Microservice architectures at companies like Netflix or Amazon are essentially SRP scaled to the cloud—each service does exactly one thing well.}
  \end{block}
\end{frame}

% ==========================================
% SLIDE 6: AI Notes
% ==========================================
\begin{frame}{6. AI Notes: Active Learning \& Geometry}
  \vspace{0.3cm}
  \textbf{Why split \texttt{xs} and \texttt{ys}?} \\
  Because the dataset lives in two geometric spaces simultaneously:
  \begin{itemize}
      \item \textbf{Feature Space (\texttt{xs}):} Independent variables. Used for clustering and measuring distance.
      \item \textbf{Objective Space (\texttt{ys}):} Dependent variables. Used to evaluate "fitness."
  \end{itemize}

  \vspace{0.4cm}
  \textbf{AI Tip: Active Learning \& Sparsity} \\
  \texttt{ezr} is an active learner. We get the absolute best results when learners are allowed to \textit{pick their own data}. By focusing only on the most informative examples (guided by our \texttt{ys}), we can dodge noise, ignore superfluous signals, and learn optimal configurations from vastly less data.
\end{frame}

